\chapter{The spherical converse: simplicity and capstone}
\label{chap:sphere_converse}

% Split from the original Gluck_Sphere.tex (stage sections preserved verbatim;
% all labels and \lean pins unchanged).

\paragraph{Formalization map.}
This chapter corresponds primarily to \texttt{Gluck/Sphere/Reconstruction.lean} and \texttt{Gluck/Sphere/Converse.lean}.

\section{Simplicity and capstone (S2-E)}
\label{sec:sphere_simplicity}

\begin{lemma}[constant curvature branch: the centered circle realizes constant $\kappa$]
  \leanok
  \label{lem:spherical_circle_realizes}
  \lean{Gluck.sphericalCircle_realizes}
  \uses{def:realizes_spherical_curvature, lem:spherical_speed_circle}
  For every constant $c>0$ the centered circle $z(\theta)=-r^\ast\,i\,e^{i\theta}$ with
  $r^\ast=\sqrt{1+c^2}-c\in(0,1)$ is a simple closed curve realizing the constant
  curvature function $\kappa\equiv c$
  (\cref{def:realizes_spherical_curvature}). This discharges the constant disjunct of
  \cref{def:four_vertex_condition} in the capstone, exactly as the Euclidean
  \texttt{gluck\_converse} handles its constant branch with a round circle; the
  reconstruction flow and the degree argument are NOT needed for it.
\end{lemma}

\begin{proof}
  \uses{lem:spherical_speed_circle}
  \leanok
  $z$ is $C^1$ (an explicit exponential), $2\pi$-periodic, and
  $z'(\theta)=r^\ast e^{i\theta}\ne0$, so the tangent-angle witness is
  $\varphi(\theta)=\theta$ with $\|z'\|=r^\ast$. Confinement: $\|z\|=r^\ast<1$ since
  $\sqrt{1+c^2}<1+c$. The speed relation is the geodesic-circle identity: with
  $\langle z,ie^{i\theta}\rangle=-r^\ast$ the right-hand side is
  $(c+r^\ast)r^\ast$ and the left-hand side $\tfrac{1+r^{\ast2}}{2}$; they agree by the
  circle identity $1+r^{\ast2}=2r^\ast(c+r^\ast)$ (defining property of
  $r^\ast=\sqrt{1+c^2}-c$; this is the content of \cref{lem:spherical_speed_circle} read
  as an identity, with $c=(1-r^{\ast2})/(2r^\ast)$). Simplicity: on $[0,2\pi)$ the map
  $\theta\mapsto -r^\ast ie^{i\theta}$ is injective because $e^{i\theta}$ is.
\end{proof}

\begin{lemma}\leanok
[spherical realization transfers under reparametrization]
  \label{lem:realizes_spherical_comp}
  \lean{Gluck.realizesSphericalCurvature_comp}
  \uses{def:realizes_spherical_curvature, lem:realizes_curvature_comp}
  If $z$ realizes the spherical curvature $\mu$
  (\cref{def:realizes_spherical_curvature}) and $\psi:\mathbb R\to\mathbb R$ is $C^1$
  with $\psi'>0$ everywhere, then $z\circ\psi$ realizes $\mu\circ\psi$.
\end{lemma}

\begin{proof}

\leanok
  Mirror of the Euclidean \cref{lem:realizes_curvature_comp}
  (\texttt{Gluck.realizesCurvature\_comp}): the chain rule gives
  $(z\circ\psi)'=\psi'\cdot(z'\circ\psi)$, so the speed scales by $\psi'>0$, the
  tangent-angle witness is $\varphi\circ\psi$, and the spherical speed relation of
  \cref{def:realizes_spherical_curvature} at $\psi(t)$ multiplies through by
  $\psi'(t)>0$ on both sides — the extra conformal factor $(1+\|z\|^2)/2$ and the
  bracket term are pointwise in the curve value and transport unchanged. Confinement and
  regularity are pointwise. This is what pulls a realization of $\kappa\circ h_1$ back
  to a realization of $\kappa$ in the capstone: apply it with $\psi=h_1^{-1}$, whose
  existence, $C^1$ regularity, positive derivative, and circle-equivariance are exactly
  the PUBLIC in-tree lemma \texttt{Gluck.exists\_C1\_circle\_inverse}
  (\texttt{Reduction.lean}, \cref{lem:exists_c1_circle_inverse}) — reuse it directly,
  no local replication needed.
\end{proof}

The simplicity of the closing trajectory reduces to the Euclidean chord-integral
machinery of \cref{chap:simplicity} through one observation: in the gauge
$\varphi(\theta)=\theta$ the disk representative is an ordinary plane curve with
$\tilde z'=\rho(\theta)\,e^{i\theta}$, $\rho>0$, i.e.\ (up to the additive constant
$\tilde z(0)$) exactly a Euclidean reconstruction curve
$\mathrm{reconstruct}\,\rho$ (\cref{def:reconstruct}) for the \emph{trajectory speed}
$\rho(\theta)=q_\kappa(\theta,\tilde z(\theta))$. The disk metric plays no role in
injectivity. The chain below works this out in three finite steps and one assembly.

\begin{lemma}\leanok
[trajectory speed: continuity, periodicity, positivity]
  \label{lem:spherical_trajectory_speed}
  \lean{Gluck.sphericalTrajectory_speed}
  \uses{def:spherical_speed, def:periodic_extension, lem:periodic_extension_periodic,
        lem:reconstruction_ode}
  Let $\kappa$ be continuous and $2\pi$-periodic (no positivity — the whole
  chain is sign-agnostic and is consumed verbatim by the mixed-sign stage 2),
  $R<1$, $\delta>0$, and let
  $z:[0,2\pi]\to\mathbb C$ be a trajectory of the $(\kappa,R,\delta)$-truncated flow
  (in the hypothesis form of \cref{lem:reconstruction_ode}) that is admissible and
  closed. Write $\tilde z$ for its periodic extension and set
  $\rho(t)=q_\kappa\bigl(t,\tilde z(t)\bigr)$
  (\cref{def:spherical_speed}). Then $\rho$ is continuous, $2\pi$-periodic, and
  $\rho(t)>0$ for all $t$.
\end{lemma}

\begin{proof}
  \uses{lem:reconstruction_ode, lem:periodic_extension_periodic}
  \leanok
  \emph{Extended admissibility.} For any $t$ reduce $u=t-2\pi\lfloor t/(2\pi)\rfloor
  \in[0,2\pi)$; admissibility at $u$ plus $2\pi$-periodicity of $\kappa$ and of
  $\theta\mapsto e^{i\theta}$ give $\|\tilde z(t)\|\le R$ and
  $\kappa(t)-\langle\tilde z(t),ie^{it}\rangle\ge\delta$ (the same reduction as inside
  the proof of \cref{lem:reconstruction_ode}).
  \emph{Continuity.} $\tilde z$ is differentiable (the exposed derivative conjunct of
  \cref{lem:reconstruction_ode}), hence continuous; $\rho$ is the quotient
  $\rho=\tfrac{1+\|\tilde z\|^2}{2(\kappa-\langle\tilde z,ie^{it}\rangle)}$ whose
  denominator is continuous and $\ge2\delta>0$ by extended admissibility, so $\rho$ is
  continuous.
  \emph{Periodicity.} $\tilde z(t+2\pi)=\tilde z(t)$
  (\cref{lem:periodic_extension_periodic}), $\kappa(t+2\pi)=\kappa(t)$, and
  $e^{i(t+2\pi)}=e^{it}$; all three inputs of $\rho$ are $2\pi$-periodic.
  \emph{Positivity.} Numerator $1+\|\tilde z\|^2\ge1>0$, denominator $\ge2\delta>0$.
\end{proof}

\begin{lemma}\leanok
[the closing trajectory is a translated reconstruction curve]
  \label{lem:spherical_trajectory_eq_reconstruct}
  \lean{Gluck.sphericalTrajectory_eq_reconstruct}
  \uses{lem:spherical_trajectory_speed, lem:reconstruction_ode, def:reconstruct,
        lem:reconstruct_deriv}
  In the setting of \cref{lem:spherical_trajectory_speed},
  \[
    \tilde z(t)\;=\;\tilde z(0)+\mathrm{reconstruct}\,\rho\,(t)
    \qquad\text{for all } t\in\mathbb R .
  \]
\end{lemma}

\begin{proof}
  \uses{lem:spherical_trajectory_speed, lem:reconstruct_deriv}
  \leanok
  Both sides are differentiable on all of $\mathbb R$ with the \emph{same} derivative
  $\rho(t)e^{it}$: the left side by the exposed derivative conjunct of
  \cref{lem:reconstruction_ode} (which is $q_\kappa(t,\tilde z(t))e^{it}
  =\rho(t)e^{it}$ by definition of $\rho$), the right side by
  \cref{lem:reconstruct_deriv} ($\rho$ is continuous by
  \cref{lem:spherical_trajectory_speed}). Their difference has vanishing derivative on
  the connected set $\mathbb R$, hence is constant; at $t=0$ the difference is
  $\tilde z(0)$ since $\mathrm{reconstruct}\,\rho\,(0)=\int_0^0=0$.
\end{proof}

\begin{lemma}\leanok
[simplicity is translation-invariant]
  \label{lem:is_simple_closed_const_add}
  \lean{Gluck.isSimpleClosed_const_add}
  \uses{def:is_simple_closed, def:is_closed_curve}
  If $\gamma$ is a simple closed curve and $w\in\mathbb C$, then $t\mapsto w+\gamma(t)$
  is a simple closed curve.
\end{lemma}

\begin{proof}

\leanok
  Closedness: $w+\gamma(t+2\pi)=w+\gamma(t)$ (\cref{def:is_closed_curve} is
  $2\pi$-periodicity, which survives adding a constant). Injectivity on $[0,2\pi)$:
  $w+\gamma(s)=w+\gamma(t)$ iff $\gamma(s)=\gamma(t)$ — cancel $w$.
\end{proof}

\begin{lemma}\leanok
[simplicity]
  \label{lem:spherical_simplicity}
  \lean{Gluck.spherical_simplicity}
  \uses{lem:spherical_trajectory_speed, lem:spherical_trajectory_eq_reconstruct,
        lem:is_simple_closed_const_add, cor:positive_curvature_implies_simple,
        def:error_vector, lem:reconstruction_ode}
  In the setting of \cref{lem:spherical_trajectory_speed} (continuous
  $2\pi$-periodic $\kappa$ — no positivity, $R<1$, $\delta>0$, an admissible
  closed trajectory $z$ of the truncated flow on
  $[0,2\pi]$), the periodic extension $\tilde z$ is a \emph{simple} closed curve:
  distinct parameters in $[0,2\pi)$ give distinct disk points. (Confinement
  $\|\tilde z\|\le R<1$ already places the curve in the open disk, so this is
  simplicity on $S^2$ via the stereographic chart; no further transport is needed.)
\end{lemma}

\begin{proof}
  \uses{lem:spherical_trajectory_eq_reconstruct, cor:positive_curvature_implies_simple,
        lem:is_simple_closed_const_add}
        \leanok
  Let $\rho$ be the trajectory speed of \cref{lem:spherical_trajectory_speed}:
  continuous, $2\pi$-periodic, strictly positive. Its error vector
  (\cref{def:error_vector}) vanishes:
  $\mathrm{errorVector}\,\rho=\mathrm{reconstruct}\,\rho\,(2\pi)
  =\tilde z(2\pi)-\tilde z(0)=0$ by \cref{lem:spherical_trajectory_eq_reconstruct} and
  $2\pi$-periodicity of $\tilde z$. Hence
  \cref{cor:positive_curvature_implies_simple} (the public Euclidean
  \texttt{isSimpleClosed\_reconstruct}; the \texttt{private}
  \texttt{simplicity\_transport} of \texttt{DahlbergStep2} is not needed) makes
  $\mathrm{reconstruct}\,\rho$ a simple closed curve, and
  \cref{lem:spherical_trajectory_eq_reconstruct} +
  \cref{lem:is_simple_closed_const_add} (translation by $\tilde z(0)$) transfer
  simplicity to $\tilde z$.
\end{proof}
